\section{Hypothetical Extended-Statistics Substrate}
\label{sec:substrate}

\subsection{Design goals}

We define the hypothetical extended-statistics substrate as a backend-local
layer that binds frozen, advisor-derived native MCV/FD payload states---or an
explicit \texttt{ABSENT\_NATIVE} state---to candidate statistics objects,
exposes an ordered active subset to PostgreSQL's statistics lookup path, and
leaves cardinality-estimator semantics inside PostgreSQL.

The substrate separates a configuration decision from the physical
statistics state that would ordinarily be created for that decision.  Its
purpose is to let the advisor compare many supported extended-statistics
configurations without installing and analyzing every design.  The
evaluation boundary therefore has five requirements: payload acquisition is
outside the design loop; no candidate design causes a per-design search-loop
\texttt{ANALYZE}; PostgreSQL's own cardinality estimator remains authoritative;
hypothetical state is confined to the advisor backend; and activation is
deterministic for a representable candidate design.

The last requirement includes a state that is easy to lose in a byte-oriented
interface.  PostgreSQL can legitimately register an extended-statistics
definition while producing no requested native payload for it.  The substrate
must preserve that \texttt{ABSENT\_NATIVE} outcome rather than silently
dropping the candidate or fabricating an empty blob.  These are engineering
requirements for the supported advisor path, not a generic DBMS
virtualization model.

The implementation is deliberately bounded.  It targets PostgreSQL 16.14,
one base relation and base-relation selection estimation, a fixed externally
supplied statistics target, arity-two MCV and functional-dependency
(``dependencies'') candidates, and the supported predicate fragment.  The
substrate does not attempt to make arbitrary PostgreSQL statistics objects or
join plans hypothetical.  Figure~\ref{fig:substrate} summarizes the boundary
between captured inputs, derived payload state, backend-local activation, and
native estimation.

\begin{figure*}[t]
\centering
\begin{tikzpicture}[
  substrate box/.style={draw, rounded corners=2pt, align=left,
    text width=3.25cm, minimum height=2.25cm, inner sep=5pt,
    font=\footnotesize},
  substrate arrow/.style={-{Latex[length=2mm]}, thick},
  substrate label/.style={font=\scriptsize, align=center, fill=white,
    inner sep=1.5pt}
]
\node[substrate box, fill=gray!12] (capture) {\textbf{Production / capture}\\[2pt]
  persistent input\\
  persisted sample\\
  workload + truth\\
  schema + compatibility};
\node[substrate box, fill=blue!10, right=8mm of capture] (repository) {\textbf{Advisor-derived repository}\\[2pt]
  $P_1,\ldots,P_n$\\
  \texttt{ABSENT\_NATIVE}\\
  one frozen realization\\
  no payload bytes in capture};
\node[substrate box, fill=orange!16, dashed, right=8mm of repository] (overlay) {\textbf{Backend-local overlay}\\[2pt]
  registered candidate state\\
  payload or absent marker\\
  active ordered subset $Y$\\[2pt]
  \scriptsize session lifetime;\\[-1pt]
  \scriptsize not persistent deployment};
\node[substrate box, fill=green!10, right=8mm of overlay] (planner) {\textbf{PostgreSQL native planner / CE}\\[2pt]
  lookup / filter\\
  deserialization\\
  MCV/FD semantics\\
  \texttt{EXPLAIN} estimate};

\draw[substrate arrow] (capture) -- node[substrate label, above]
  {read-only\\reconstruction} (repository);
\draw[substrate arrow] (repository) -- node[substrate label, above]
  {register\\candidate state} (overlay);
\draw[substrate arrow] (overlay) -- node[substrate label, above]
  {native\\lookup} (planner);
\node[font=\bfseries\scriptsize, align=center, below=6mm of overlay.south east,
  anchor=north east] (invariant) {statistics state is virtualized;\\
  estimator semantics stay native};
\end{tikzpicture}
\caption{The substrate derives native candidate state from a persisted capture,
registers it in a backend-local overlay, and exposes an ordered active subset
to PostgreSQL's native statistics path.  The overlay virtualizes state; the
planner remains responsible for estimator semantics.}
\label{fig:substrate}
\end{figure*}

\subsection{Candidate payload repository}

The system keeps four related objects distinct.  A \emph{candidate
definition} is the deterministic catalog entry: relation, attributes,
mechanism, statistics-object definition, and precedence rank.  A \emph{shell}
is the corresponding PostgreSQL extended-statistics object used to give the
candidate a catalog identity in the private advisor database.  A \emph{native
payload} is the serialized value PostgreSQL produced for that shell under one
statistics realization.  The \emph{payload repository} binds definitions to
those payload states and records acquisition provenance, integrity digests,
and PostgreSQL build identity.

The repository is derived from a frozen realization, not supplied by the
production capture contract.  Production capture exports the persisted sample
and the workload, truth, schema, and compatibility metadata needed to
reconstruct it.  It explicitly does not contain ordinary \texttt{pg\_statistic}
rows, extended-statistics payload bytes, or a derived repository.  In the
private advisor, the sample is staged, the frozen-sample controls are applied,
and PostgreSQL derives ordinary and extended statistics from that sample.  On
a repository cache miss, the acquisition path creates candidate shells,
applies the fixed target, runs one relation-grouped \texttt{ANALYZE}, and
reads the native MCV or dependency value through PostgreSQL's send function.
On a cache hit, the advisor recreates the catalog shells and performs the
required pre-search shell realization before reusing the validated payload
repository.  Both paths finish before a design is evaluated.

Each repository entry has an explicit realization state.  \texttt{PRESENT}
means that the requested native payload exists and is retained as non-empty
serialized bytes with a digest.  \texttt{ABSENT\_NATIVE} means that the
candidate definition was registered and checked, but PostgreSQL legitimately
returned a null requested payload field; the entry therefore has no payload
bytes.  \texttt{UNREGISTERED} is different: it denotes an acquisition or
validation failure, such as a missing catalog definition or incompatible
identity, and is never treated as a statistical outcome.  The adapter fails
closed on that state.

Maintenance cost is attached to the candidate definition.  It is therefore
the same design-level property whether the frozen realization records
\texttt{PRESENT} or \texttt{ABSENT\_NATIVE}; payload presence is not used as a
proxy for cost.  The maintenance model and budget are described separately in
Section~6.

\subsection{Backend-local registration}

Registration is performed by the one-session \texttt{PostgresAdapter} after
the repository has been loaded.  The adapter first resets the overlay with
\texttt{pg\_hypothetical\_extstats\_reset()} and resolves each candidate's
statistics-object identity.  It checks the relation in \texttt{pg\_class},
looks up the candidate shell in \texttt{pg\_statistic\_ext} (falling back to
the deterministic shell name when the stored backend OID has changed), and
verifies both relation ownership and the requested mechanism kind.

For a \texttt{PRESENT} entry it calls the function
\texttt{pg\_hypothetical\allowbreak\_extstats\allowbreak\_register} with
\texttt{(statoid, relid, kind, payload)}.
The patched function validates non-null OIDs, checks the shell's actual
relation, validates the mechanism, and invokes PostgreSQL's native MCV or
dependency deserializer before copying the bytes into backend-local state.
For \texttt{ABSENT\_NATIVE}, the adapter calls
\texttt{pg\_hypothetical\allowbreak\_extstats\allowbreak\_register\_absent}
with \texttt{(statoid, relid, kind)};
the patch records the absence without inventing a serialized value.  A
duplicate registration, unknown shell, relation mismatch, kind mismatch, or
malformed payload is an error rather than an implicit fallback.

The patch stores the repository in memory contexts rooted at the current
backend's top memory context and stores the active design in a child context.
Registration and activation consequently live for the backend session unless
the adapter resets them.  Registration does not insert or update persistent
catalog rows.  The shell definitions used for identity are created only in
the private advisor database during acquisition or shell reconstruction; the
registration calls themselves perform catalog checks and populate the overlay.
Production PostgreSQL is never modified by this interface.

\subsection{Ordered activation}

Let $S$ be the frozen candidate catalog and let $Y\subseteq S$ be a
representable supported design.  The candidate catalog assigns each candidate
a unique precedence rank, and the catalog normalizes every design into that
fixed order.  The adapter translates the resulting candidate IDs to the
resolved shell OIDs and calls
\texttt{pg\_hypothetical\_extstats\_activate(oid[])}.  The order of this
array is retained by the backend and is the effective order supplied to the
native statistics lookup path.

Activation validates the whole replacement before publishing it: every OID
must have a registered payload state, no OID may occur twice, and the
relation identity recorded at registration is retained.  The backend builds
an active-membership hash table and an ordered OID array atomically, then
replaces the previous active context.  The adapter immediately reads
\texttt{pg\_hypothetical\_extstats\_active()} and rejects a mismatch with
the requested order.

When the planner asks for extended statistics on the targeted relation, the
patched filter first recognizes that a hypothetical realization is active,
then emits the active registered OIDs in the requested order.  Registered
candidate shells not in $Y$ are omitted from that relation's hypothetical
list; unrelated, unregistered catalog objects remain ordinary PostgreSQL
objects.  Thus the interface represents configurations over the frozen
candidate catalog and its supported ordering, not every mathematical subset
of every PostgreSQL statistics object.

\subsection{Native planner hooks}

The central invariant is that the overlay virtualizes statistics state; it
does not reimplement cardinality-estimator semantics.  The patch adds a
narrow hook after PostgreSQL obtains a relation's extended-statistics OID list:
the list is filtered and reordered before the existing planner code creates
its \texttt{StatisticExtInfo} nodes.  If a registered active entry is marked
\texttt{ABSENT\_NATIVE}, the corresponding native loader is allowed to return
no value.  If it is \texttt{PRESENT}, the MCV or dependency loader receives the
registered bytes and invokes PostgreSQL's existing deserializer.

All subsequent supported behavior remains in PostgreSQL.  The planner applies
statistics applicability and lookup, determines effective precedence and
selection, consumes clauses, deserializes native MCV or dependency state,
combines MCV and FD information, and computes the numerical selectivity and
cardinality estimate.  The adapter does not calculate an MCV match, an FD
degree, a residual clause state, or a replacement estimate in Python.  The
advisor supplies hypothetical statistics state; PostgreSQL executes the
estimator semantics that consume it.

This boundary also explains why the repository is more than a byte cache.  A
cache can retain a serialized value, but the substrate additionally validates
candidate identity and relation ownership, preserves a legitimate native
absence, controls active membership and order, and routes the registered state
through the native lookup and loading path.  These behaviors make a stored
payload usable as a PostgreSQL what-if state within the supported fragment.

\subsection{Isolation and state lifetime}

The hypothetical repository and active ordered subset are backend-local.  A
reset discards both registration and activation state for that advisor
session; no persistent hypothetical statistics object is installed in a
production catalog.  The private patched PostgreSQL is therefore an
evaluation environment, not a replacement production server.  Once advice is
accepted, the selected design is rendered as ordinary PostgreSQL
\texttt{CREATE STATISTICS}/\texttt{ALTER STATISTICS} DDL for DBA review and
later physical collection.  The distinction is intentional: registering a
payload in the overlay changes what the private planner sees, whereas
physical DDL followed by \texttt{ANALYZE} changes persistent production
statistics state.

\subsection{Supported substrate boundary}

Table~\ref{tab:substrate-scope} fixes the boundary for the claims in this
paper; extensions outside these rows are a new scope.

\begin{table}[t]
\caption{Supported hypothetical extended-statistics substrate.}
\label{tab:substrate-scope}
\centering
\small
\begin{tabular}{p{0.29\columnwidth}p{0.62\columnwidth}}
\toprule
Dimension & Supported boundary \\
\midrule
PostgreSQL & 16.14 advisor build and patch \\
Relation scope & one base relation; base-relation selection CE \\
Mechanisms & arity-two MCV and dependencies (FD) \\
Target & externally supplied fixed target, default $T=100$ \\
Realization & one persisted frozen sample and derived repository \\
Workload & supported single-relation predicate forms \\
Excluded & joins, higher arity, other mechanisms, target optimization \\
Portability & no cross-version claim \\
\bottomrule
\end{tabular}
\end{table}

The table is a scope contract, not a claim that the substrate enumerates all
possible configurations within PostgreSQL.  In particular, same-realization
fidelity below does not imply that a future native \texttt{ANALYZE} will
produce identical payload bytes or estimates.
